% --------------------------- PRIMITIVES ----------------------------
\newcommand{\insmall}[1]{\text{\small{#1}}}
%\newcommand{\minus}{\insmall{-}}
% Use \AtBeginDocument to ensure unicode-math doesn't overwrite it

\AtBeginDocument{
  \renewcommand{\minus}{%
    \mathbin{%
      \mathchoice
        {\scalebox{0.8}{$\text{\bf -}$}} % Display and Text style
        {\scalebox{0.8}{$\text{\bf -}$}} % Script style
        {\scalebox{0.5}{$\text{\bf -}$}} % Script-script style
        {\scalebox{0.8}{$\text{\bf -}$}} % Default fallback
    }%
  }
}

% The field C and the usual subsets R, Q, Z, N.
% NOTE: \mathbf produces bold italic in unicode-math, switch to \mathbb for standard notation
\newcommand\usualSet[1]{\mathbf{#1}} 
\def\C{\usualSet{C}}
\def\R{\usualSet{R}}
\def\Q{\usualSet{Q}}
\def\Z{\usualSet{Z}}
\def\N{\usualSet{N}}
\def\K{\usualSet{K}}

\def\Def{\triangleq}
\def\defIFF{\quad \text{\bf iff} \quad} % Halmos' iff

\newcommand\counting[1]{#1=1, 2, 3, \dots} 
\newcommand\naturals[1]{#1=0, 1, 2, \dots} 

\newcommand{\up}[1]{^{(#1)}}
\newcommand{\low}[1]{_{#1}}
\newcommand{\upnw}[2]{\fourIdx{#2}{}{}{}#1}
\newcommand{\downsw}[2]{\fourIdx{}{#2}{}{}#1}
\newcommand{\scriptatleft}[3]{\fourIdx{#2}{#3}{}{}#1}
\newcommand{\diagscript}[3]{\fourIdx{#2}{}{}{#3}#1}

\newcommand{\function}[1]{\mathtt{#1}}
\newcommand{\relation}[2]{{#1}_{#2}}
\newcommand{\f}[2]{#1(#2)}
\newcommand{\id}[1]{\text{id}_{#1}}
\renewcommand{\restriction}[2]{#1\vert_{#2}}

% Predicate calculus, logic
% % Bunch of synonyms
\def\then{\Rightarrow} 
\def\therefore{\Rightarrow}
\def\implies{\Rightarrow}
\def\since{\Leftarrow}
\def\because{\Leftarrow} 
\renewcommand{\iff}{\Leftrightarrow} 

% % Iverson bracket notation
%\newcommand{\Iverson*[\big]}[2][\big]{#1[\, #2\, #1]}
%\newcommand{\Iverson*[\big]}[2][\big]{#1[\, #2\, #1]}
\NewDocumentCommand{\Iverson}{s o m}{%
  % #1: * (optional star)
  % #2: [\big] (optional size)
  % #3: Content (required)
  \IfBooleanTF{#1}{
    \left[ \hspaceNorm #3 \hspaceNorm \right]
  }{
    \IfValueTF{#2}{
      #2[ \hspaceNorm #3 \hspaceNorm #2]
    }{
      \errmessage{Iverson: Size argument needed.}
    }
  }
}

% Sets
% % Braces
\DeclarePairedDelimiter\singleton\{\}
\newcommand{\set}[3][\big]{\singleton[#1]{#2: #3}}

% % Union, intersection. 
\def\cuts{\cap} % synonym 

% % Subset, supset
\AtBeginDocument{\let\subset\subseteq
\let\supset\supseteq}
\def\contains{\supseteq} % synonym

% Vector spaces
% % Convex hull
\DeclareMathOperator\co{co} 

% % Span
\DeclareMathOperator\vect{span}

% % Misc
\newcommand{\bigsetminus}[1][\big]{\mathbin{#1\backslash}}
\DeclareMathOperator\opcard{card} 

% Topology 
\newcommand{\interior}[1]{{#1}^\circ}

\AtBeginDocument{
  \DeclareRobustCommand{\closure}[1]{\overline{#1}}
}

\newcommand{\localbase}[1]{\mathcal #1}
\def\weakstar{\text{weak}^\ast\text{-}}

% Standard calculus
% % Intervals, as partial functions of the Interval command below.
% % % Open.
\NewDocumentCommand{\openInterval}{s o m m}{%
  \IfBooleanTF{#1}{
    \Interval*{#3}{#4}{]}{[}
  }{% 
    \IfValueTF{#2}{% If [size]:
      \Interval[#2]{#3}{#4}{]}{[}
    }{% 
      % If no [size], use standard delimiters with correct spacing
      \errmessage{openInterval: Size argument is required.}
    }
  }
}

% % % Closed.
\NewDocumentCommand{\closedInterval}{s o m m}{%
  \IfBooleanTF{#1}{
    \Interval*{#3}{#4}{[}{]}
  }{% 
    \IfValueTF{#2}{% If [size]:
      \Interval[#2]{#3}{#4}{[}{]}
    }{% 
      % If no [size], use standard delimiters with correct spacing
      \errmessage{closedInterval: Size argument is required.}
    }
  }
}

% % % Open-closed ]...].
\NewDocumentCommand{\openClosedInterval}{s o m m}{%
  \IfBooleanTF{#1}{
    \Interval*{#3}{#4}{]}{]}
  }{% 
    \IfValueTF{#2}{% If [size]:
      \Interval[#2]{#3}{#4}{]}{]}
    }{% 
      % If no [size], use standard delimiters with correct spacing
      \errmessage{openClosedInterval: Size argument is required.}
    }
  }
}

% % % Closed-open [...[.
\NewDocumentCommand{\closedOpenInterval}{s o m m}{%
  \IfBooleanTF{#1}{
    \Interval*{#3}{#4}{[}{[}
  }{% 
    \IfValueTF{#2}{% If [size]:
      \Interval[#2]{#3}{#4}{[}{[}
    }{% 
      % If no [size], use standard delimiters with correct spacing
      \errmessage{closedOpenInterval: Size argument is required.}
    }
  }
}

% % % Interval, the primitive. 
% % % Should not be called directly. 
% % % Use open Interval, ..., closedOpenInterval; see above
\NewDocumentCommand{\Interval}{s o m m m m}{%
  \IfBooleanTF{#1}{% Auto-sizing
    \left#5 \hspaceNorm #3, #4 \hspaceNorm  \right#6
  }{% Else, manual sizing
    \IfValueTF{#2}{% If [size] 
      \mathopen{#2#5} \hspaceNorm  #3, #4 \hspaceNorm  \mathclose{#2#6}
    }{% Else, error
      \errmessage{Interval: Size argument is required.}
    }
  }
}

%newcommand{\set}[3][\big]{\singleton[#1]{#2: #3}}

% % Usual spaces
\newcommand{\Continuous}{{\mathcal C}} 
\newcommand{\D}{{\mathcal D}} 

% % Bounds
\newcommand{\supp}[1]{\operatorname{supp}(#1)}

% % Usual functions and operator
\DeclareMathOperator{\cotan}{cotan}
\DeclareMathOperator\sgn{sgn}
\NewDocumentCommand{\ceil}{s o m}{%
  \IfBooleanTF{#1}{%
    % Star version: auto-sizing with \left\right
    \IfValueTF{#2}{
      \errmessage{ceil: Star mode. Don't specify the [size].}
    }{
      \lceil \hspaceNorm #3 \hspaceNorm \rceil
    }
  % Non-star version: manual sizing
  }{
    \IfValueTF{#2}{% if [size] 
      #2\lceil \hspaceNorm #3 \hspaceNorm #2\rceil
    }{% elif no [size]
      \errmessage{ceil: Size command required.}
    }
  }
}

\newcommand{\floor}[1]{\lfloor \hspaceNorm #1 \hspaceNorm\rfloor}
\DeclareMathOperator\Leibnizdiff{d}
\newcommand{\diff}[1]{\Leibnizdiff{\!#1}}

% % Evaluation
\newcommand{\evalbar}[4][\Big]{#2\mathrel{#1|}_{#3}^{#4}}
\newcommand{\evalbracket}[4][\Big]{#1[\,#2\,#1]_{#3}^{#4}}

% % Limits
\newcommand{\tendsto}[2]{\xrightarrow[#1 \to #2]{}}

% % Norms
\def\hspaceNorm{\hspace{0.12 em}}

% % Alternative to \abs
\NewDocumentCommand{\magnitude}{s o m}{%
  \IfBooleanTF{#1}{%
    % Star version: auto-sizing with \left\right
    \IfValueTF{#2}{
      \errmessage{magnitude: Star mode. Don't specify the [size].}
    }{
        \left\lvert \hspaceNorm #3 \hspaceNorm #1\right\rvert
    }
  % Non-star version: manual sizing
  }{
    \IfValueTF{#2}{% if [size] 
        #2\lvert \hspaceNorm #3 \hspaceNorm #2\rvert
    }{% elif no [size]
      \errmessage{magnitude: Size command required.}
    }
  }
}
%\newcommand{\norm}[2][]{
%  \if\relax\detokenize{#1}\relax 
%    \lVert \hspaceNorm  #2 \hspaceNorm \rVert 
%  \else 
%    \left\lVert \hspaceNorm #2 \hspaceNorm \right\rVert_{#1} 
%  \fi
%}

% \norm command
% #1: * (optional star for auto-sizing)
% #2: [\big] (optional size command like \big, \Big, \bigg, \Bigg)
% #3: Subscript (optional)
% #4: Content (required)
\RenewDocumentCommand{\norm}{s o o g}{%
  \IfBooleanTF{#1}{%
    % Star version: auto-sizing with \left\right
    \IfValueTF{#2}{% if [subscript]
      \IfValueTF{#3}{ % if [content]
      \errmessage{norm: Star mode. Don't specify the [size].}%
      }{ % elif {content}
        \IfValueTF{#4}{
          {\left\lVert \hspaceNorm #4 \hspaceNorm \right\rVert}_{#2}%
        }{ % else:
        \errmessage{norm: Star mode. Don't specify the [size].}%
        }
      }
    }{% elif no [subscript]
      \left\lVert \hspaceNorm #4 \hspaceNorm \right\rVert
    }%
  }{%
    % Non-star version: manual sizing
    \IfValueTF{#2}{% if [size]
      \IfValueTF{#3}{% if [p] 
        {#2\lVert \hspaceNorm #4 \hspaceNorm #2\rVert}_{#3}%
      }{% elif no [p]
        % We assume that #2 was a legal [size] value
        #2\lVert \hspaceNorm #4 \hspaceNorm #2\rVert
      }%
    }{% elif no [size]
      \errmessage{norm: Size command required. Use \string\norm[\string\big]{F} or \string\norm*{F}}%
    }
  }
}


% % the ||| . ||| version
\newcommand{\vertiii}[1]{{\left\vert\kern-0.25ex\left\vert\kern-0.25ex\left\vert #1 
    \right\vert\kern-0.25ex\right\vert\kern-0.25ex\right\vert}}
\newcommand{\normtriple}[2]{\vertiii{\hspaceNorm #2 \hspaceNorm}_{#1}}

% Typesetting

\newcommand{\varit}[1]{\mathit{#1}}
\def\vart{\varit{t}}

\def\ie{\textit{i.e., }}
\def\eg{\textit{e.g., }}
\def\cf{\textit{cf.\,}} 
\def\IFF{{\bf iff }}
\def\wlg{{\bf wlg }}